% BHU PYQ -- Manually transcribed & error-corrected
% Paper: BCH-226 : Business Statistics
% Exam: B.Com. (Hons.) (Semester IV) Examination, 2015-16

\documentclass[11pt,a4paper]{article}
\usepackage[a4paper,top=2.5cm,bottom=2.5cm,left=2.8cm,right=2.8cm,headheight=14pt]{geometry}
\usepackage{amsmath,amssymb}
\usepackage[T1]{fontenc}
\usepackage[utf8]{inputenc}
\usepackage{lmodern,microtype,fancyhdr,enumitem,hyperref,lastpage,parskip}

\hypersetup{
    colorlinks=true,
    urlcolor=black,
    linkcolor=black,
    pdftitle={BCH-226: Business Statistics -- BHU B.Com. (Hons.) Sem IV 2015-16}
}

\pagestyle{fancy}
\fancyhf{}
\renewcommand{\headrulewidth}{0.4pt}
\renewcommand{\footrulewidth}{0.4pt}
\fancyhead[L]{\small\textit{Banaras Hindu University}}
\fancyhead[R]{\small\textit{BCH-226: Business Statistics}}
\fancyfoot[C]{\small Page \thepage\ of \pageref{LastPage}}

\newlist{parts}{enumerate}{2}
\setlist[parts,1]{label=(\alph*),leftmargin=*,itemsep=5pt,topsep=3pt}
\setlist[parts,2]{label=(\roman*),leftmargin=*,itemsep=3pt,topsep=2pt}
\newcommand{\pts}[1]{\hfill\textit{[#1]}}

\begin{document}

\begin{center}
  {\large\bfseries BANARAS HINDU UNIVERSITY}\\[2pt]
  {\normalsize B.Com. (Hons.) --- Semester IV Examination, 2015-16}\\[6pt]
  \rule{0.8\linewidth}{0.5pt}\\[6pt]
  {\large\bfseries Paper No. BCH-226: Business Statistics}\\[4pt]
  {\normalsize Time: 3 Hours\hfill Maximum Marks: 70}
\end{center}

\rule{\linewidth}{0.4pt}

\smallskip
\noindent\textit{Instructions: Answer all questions. All questions carry equal marks (14 marks each).}

\bigskip

\noindent\textbf{Question 1.} \\
What are Index Numbers? How are they constructed? Explain the role of weights in the construction of a general price index number.\pts{14}

\centerline{\textbf{OR}}

The following table gives the average wholesale prices of three commodities for the years from 2011 to 2015. Find Chain Base Index Numbers with 2011 as the base year:

\begin{center}
\small
\begin{tabular}{|c|c|c|c|c|c|}
\hline
\textbf{Items} & \textbf{2011} & \textbf{2012} & \textbf{2013} & \textbf{2014} & \textbf{2015} \\
\hline
I & 2 & 3 & 4 & 5 & 6 \\
II & 8 & 10 & 12 & 15 & 18 \\
III & 4 & 5 & 8 & 10 & 12 \\
\hline
\end{tabular}
\end{center}\pts{14}

\medskip\rule{\linewidth}{0.2pt}

\pagebreak

\noindent\textbf{Question 2.} \\
From the following data, calculate Fisher's Ideal Index Number. Does it satisfy the Factor Reversal and Time Reversal Tests?

\begin{center}
\small
\begin{tabular}{|c|c|c||c|c|}
\hline
\textbf{Commodity} & \multicolumn{2}{c||}{\textbf{2014 Base Year}} & \multicolumn{2}{c|}{\textbf{2015 Current Year}} \\
\cline{2-5}
& \textbf{Price (Rs.)} & \textbf{Quantity} & \textbf{Price (Rs.)} & \textbf{Quantity} \\
\hline
A & 6 & 50 & 10 & 56 \\
B & 2 & 100 & 2 & 120 \\
C & 4 & 60 & 6 & 60 \\
D & 10 & 30 & 12 & 24 \\
E & 8 & 40 & 12 & 36 \\
\hline
\end{tabular}
\end{center}\pts{14}

\centerline{\textbf{OR}}

An enquiry into the budget of middle class families in Allahabad gave the following information:

\begin{center}
\small
\begin{tabular}{|l|c|c|c|c|c|}
\hline
\textbf{Item} & \textbf{Food} & \textbf{Rent} & \textbf{Clothing} & \textbf{Fuel} & \textbf{Miscellaneous} \\
\hline
\textbf{Expenses on} & 35\% & 15\% & 20\% & 10\% & 20\% \\
\textbf{Prices (2009) in Rs.} & 150 & 30 & 75 & 25 & 40 \\
\textbf{Prices (2015) in Rs.} & 145 & 30 & 65 & 23 & 45 \\
\hline
\end{tabular}
\end{center}

What changes in cost of living figures of 2015 as compared with that of 2009 are seen?\pts{14}

\medskip\rule{\linewidth}{0.2pt}

\pagebreak

\noindent\textbf{Question 3.} \\
What do you mean by 'secular trend'? Discuss any two methods of isolating trend values in a time-series and point out their shortcomings.\pts{14}

\centerline{\textbf{OR}}

Below are given the figures of production (in thousand tonnes) of a sugar factory:

\begin{center}
\small
\begin{tabular}{|l|c|c|c|c|c|c|c|}
\hline
\textbf{Year} & 2009 & 2010 & 2011 & 2012 & 2013 & 2014 & 2015 \\
\hline
\textbf{Production} & 77 & 88 & 94 & 85 & 91 & 98 & 90 \\
\hline
\end{tabular}
\end{center}

\begin{parts}
  \item Fit a straight line by the least squares method and calculate the trend values.\pts{6}
  \item What is the monthly increase in the production of sugar?\pts{4}
  \item Also estimate the production for 2018.\pts{4}
\end{parts}

\medskip\rule{\linewidth}{0.2pt}

\noindent\textbf{Question 4.} \\
Calculate the seasonal variations using the data given below:

\begin{center}
\small
\begin{tabular}{|c|c|c|c|c|}
\hline
\textbf{Year} & \textbf{Summer} & \textbf{Monsoon} & \textbf{Autumn} & \textbf{Winter} \\
\hline
2011 & 30 & 81 & 62 & 119 \\
2012 & 33 & 104 & 86 & 171 \\
2013 & 42 & 153 & 99 & 221 \\
2014 & 56 & 172 & 129 & 235 \\
2015 & 67 & 201 & 136 & 302 \\
\hline
\end{tabular}
\end{center}\pts{14}

\centerline{\textbf{OR}}

Calculate the seasonal index from the following data using the average method:

\begin{center}
\small
\begin{tabular}{|c|c|c|c|c|}
\hline
\textbf{Year} & \multicolumn{4}{c|}{\textbf{Quarterly}} \\
\cline{2-5}
& \textbf{I} & \textbf{II} & \textbf{III} & \textbf{IV} \\
\hline
2011 & 72 & 68 & 80 & 70 \\
2012 & 76 & 70 & 82 & 74 \\
2013 & 74 & 66 & 84 & 80 \\
2014 & 76 & 74 & 84 & 78 \\
2015 & 78 & 74 & 86 & 82 \\
\hline
\end{tabular}
\end{center}\pts{14}

\medskip\rule{\linewidth}{0.2pt}

\pagebreak

\noindent\textbf{Question 5.} \\
What do you understand by statistical quality control? Discuss its aspects and importance in industry.\pts{14}

\centerline{\textbf{OR}}

Explain the use of $p$-chart, $np$-chart and $c$-chart in this connection for statistical quality control.\pts{14}

\vfill
\rule{\linewidth}{0.4pt}
\begin{center}\small
\href{https://bhu-exam.vercel.app/}{Click here to visit our website}\\[4pt]
\href{https://me-aryan.vercel.app/}{Click here to visit our contributor}\\[4pt]
\href{https://quashan-me.vercel.app/}{Click here to visit our contributor}
\end{center}

\end{document}
